% Generated draft: requires mathematical and source review.
\documentclass[11pt]{article}
\usepackage[margin=25mm]{geometry}
\usepackage{fontspec}
\setmainfont{Latin Modern Roman}
\usepackage{amsmath,amssymb,mathtools,unicode-math}
\setmathfont{Latin Modern Math}
\usepackage{enumitem,needspace,xurl,hyperref}
\setlength{\parindent}{0pt}
\setlength{\parskip}{5pt}
\setlength{\emergencystretch}{3em}
\newcommand{\N}{\mathbb{N}}
\newcommand{\Z}{\mathbb{Z}}
\newcommand{\Q}{\mathbb{Q}}
\newcommand{\R}{\mathbb{R}}
\newcommand{\C}{\mathbb{C}}
\newcommand{\F}{\mathbb{F}}
\newcommand{\A}{\mathbb{A}}
\newcommand{\PP}{\mathbb{P}}
\newcommand{\E}{\mathbb{E}}
\newcommand{\eps}{\varepsilon}
\newcommand{\dd}{\,\mathrm{d}}
\begin{document}
\section*{UnsolvedMath ID: 20002086}
\textbf{Source problem code:} catalogue entry
\subsection{Definitions and mathematical statement}
\textbf{Question 1.28, discussed by Peter Ozsv\'ath.} Can the differentials of the Khovanov-to-Heegaard-Floer spectral sequence for a link in $S^3$ be computed combinatorially? Can Khovanov theory yield bounds on the smooth four-ball genus? The source motivates these as part of computing Ozsv\'ath--Szab\'o invariants combinatorially or axiomatically; its comparative comments about ease and usefulness are historical motivation, not mathematical hypotheses.

Write $\Sigma_2(K)$ for the double branched cover. A standard convention for the spectral sequence is
\[
 E_2=\widetilde{\operatorname{Kh}}(\overline K;\F_2)
       \Longrightarrow\widehat{HF}(\Sigma_2(K);\F_2),
\]
using reduced Khovanov homology and a mirror $\overline K$; equivalent orientation/mirror conventions exist. The source suppresses reduction and mirror conventions. The following rank statements do not depend on that mirror choice, but an actual differential computation does.

\textbf{Supplied distinct affirmative answers.} The aid reports that the Lipshitz--Ozsv\'ath--Thurston bordered construction computes combinatorially a filtered complex equivalent to the original Ozsv\'ath--Szab\'o complex, hence its spectral sequence. This is stronger than merely computing its endpoint but is not an assertion of an efficient closed formula for every page. Separately, the Lee--Rasmussen deformation gives, for a knot,
\[
 |s(K)|\le2g_4(K),
\]
where $g_4$ is the minimum genus of a smooth connected oriented surface in $B^4$ with boundary $K$. The usual Lee--Rasmussen construction uses characteristic-zero coefficients; it is not identified with the mod-two branched-cover spectral sequence or its page differentials. The source starts with links; the knot genus/rank specialization is explicit here.

\textbf{Supplied determinant rank budget for knots.} For a knot $K$, let
\[
 R=\dim_{\F_2}\widetilde{\operatorname{Kh}}(K;\F_2),\qquad
 F=\dim_{\F_2}\widehat{HF}(\Sigma_2(K);\F_2),\qquad
 D=\det K=|H_1(\Sigma_2(K);\Z)|.
\]
Then $D>0$ and the aid uses $R\ge F\ge D$, with matching parity, to define nonnegative integers
\[
 e(K)=\frac{R-D}{2},\qquad h(K)=\frac{F-D}{2}.
\]
For the bounded finite-dimensional spectral sequence, let $\operatorname{rank}(d_r)$ mean total rank on the whole page. The exact telescoping identity supplied by the aid is
\[
 \sum_{r\ge2}\operatorname{rank}(d_r)
     =\frac{R-F}{2}=e(K)-h(K).
\]
It is a necessary rank test for a proposed filtered model, not a determination of the differentials or filtered extension data. If $e=1$, the only rank-level possibilities are $h=1$ with no higher differential, or $h=0$ with total higher-differential rank one. If $e=0$, all higher differentials vanish.

For positive torus knots $K=T(2,2m+1)$, $m\ge0$, the aid gives $R=F=D=2m+1$, hence collapse from $E_2$, while $g_4(K)=m$ is unbounded. Thus collapse or zero determinant budget does not imply small four-ball genus. The genus bound from the Lee--Rasmussen deformation and the rank budget answer different aspects of the source questions; no classification of all page behavior or link-genus conventions is supplied.
\subsection{Short English statement}
The aid reports separate combinatorial answers for the branched-cover spectral sequence and Rasmussen's four-ball-genus bound. Its extra rank identity counts total spectral cancellations using Khovanov rank, Floer rank and determinant. The collapsing family $T(2,2m+1)$ has unbounded genus, so collapse alone does not give a genus bound.
\subsection{Sources}
\begin{itemize}
\item Local UnsolvedMath registry, permanent ID 20002086, source code catalogue entry; canonical \texttt{aim-geometry-notes.json}, notes[423].
\item AIM \emph{Holomorphic curves in contact geometry}, Question 1.28, discussion by Peter Ozsv\'ath: \url{https://aimath.org/WWN/contactgeom2/contactgeom2.pdf}.
\item Supplied named constructions: Ozsv\'ath--Szab\'o spectral sequence, Lipshitz--Ozsv\'ath--Thurston bordered construction, Lee--Rasmussen deformation. No full titles, dates, theorem labels or proving citations are supplied in this record.
\item Registry research attempt 1, dated 2026-08-14: partial result, problem status at run solved, rigor ``proved'', novelty confidence low. The explicitly machine-generated, non-peer-reviewed aid requires independent verification before citation.
\end{itemize}
\end{document}
